%%==============================================================
%% MODELO DE TRABALHO DE CONCLUSÃO DE CURSO (VAZIO)
%% Curso de Sistemas de Informação - Universidade Federal de Viçosa
%% Campus de Rio Paranaíba
%% Baseado no Manual de Normalização de Trabalhos Acadêmicos da UFV
%% (edição 2026-2; NBR 14724/2025, 6023/2025, 10520/2023)
%%
%% Arquivo em formato UTF-8. Compilar com pdflatex + biber (nesta pasta: latexmk).
%% Precisa do arquivo abntex2-UFV.sty (pasta pai; se copiar este modelo para outro
%% lugar, copie o .sty para a mesma pasta e apague a linha ensure_path do latexmkrc).
%%
%% COMO USAR
%%  1. Preencha os dados do trabalho (seção "DADOS DO TRABALHO", logo abaixo).
%%  2. Escreva o texto nos arquivos da pasta capitulos/ (um por seção primária).
%%  3. Cadastre as referências em referencias.bib e cite com \cite{chave} e
%%     \textcite{chave}.
%%  4. Descomente as listas, pacotes e elementos opcionais de que precisar.
%%
%% Todos os recursos, com exemplos, dicas de metodologia científica e de uso da
%% classe, estão documentados no guia da pasta ../docs (guia-abntex2-ufv.pdf).
%%==============================================================

% A classe abntex2 carrega o babel com o nome obsoleto "brazil"; o aviso é inofensivo
\RequirePackage{silence}
\WarningFilter{babel}{Name 'brazil' is deprecated}
\documentclass[
    12pt,                   % tamanho da fonte
    openright,              % seções primárias começam em página ímpar (em oneside equivale a nova página)
    oneside,                % impressão só no anverso (twoside para frente e verso)
    a4paper,                % tamanho do papel
    english,                % idioma adicional (hifenização)
    brazilian,              % idioma principal do documento
    ]{abntex2}

% ---
% PACOTES
% ---
\usepackage{abntex2-UFV}        % Personalização UFV (opções: fonte=times|arial|cm, capacampus)
\usepackage[T1]{fontenc}        % Seleção de códigos de fonte de saída
\usepackage{indentfirst}        % Indenta o primeiro parágrafo de cada seção
\usepackage{graphicx}           % Inclusão de figuras
\usepackage{csquotes}           % Aspas (\enquote) e citações no padrão do idioma
\usepackage[final]{microtype}   % Melhorias de justificação

% Tabelas e quadros
\usepackage{booktabs}           % \toprule, \midrule e \bottomrule
\usepackage{tabularx}           % colunas X de largura automática
\usepackage{multirow}           % células mescladas em linhas
%\usepackage{longtable}         % tabelas que ocupam mais de uma página
%\usepackage{pdflscape}         % página em paisagem (ambiente landscape)
%\usepackage{colortbl}          % células sombreadas (\cellcolor)
%\usepackage{siunitx}           % alinhamento decimal e unidades
%\sisetup{output-decimal-marker = {,}, group-separator = {.}}

% Diagramas, gráficos e cronograma (descomente o que usar)
%\usepackage{tikz}
%\usetikzlibrary{arrows.meta, positioning}
%\usepackage{pgfplots}
%\pgfplotsset{compat=newest}

% Código-fonte (use dentro dos ambientes codigo e ilustracao). Para ativar, descomente o
% bloco inteiro; as opções estão explicadas no Apêndice D do guia (docs/guia-abntex2-ufv.pdf).
%\usepackage[scaled=0.88]{FiraMono}   % fonte de largura fixa legível
%\usepackage{xcolor}
%\usepackage{listings}
%\definecolor{ufvcodebg}{gray}{0.96}
%\definecolor{ufvkw}{RGB}{0,51,153}      % palavras-chave
%\definecolor{ufvcom}{RGB}{0,110,50}     % comentários
%\definecolor{ufvstr}{RGB}{163,21,21}    % cadeias de caracteres
%\lstdefinestyle{ufvcodigo}{
%  basicstyle=\ttfamily\footnotesize, columns=fixed, keepspaces=true, tabsize=4,
%  keywordstyle=\bfseries\color{ufvkw}, commentstyle=\itshape\color{ufvcom},
%  stringstyle=\color{ufvstr}, showstringspaces=false,
%  numbers=left, numberstyle=\tiny\color{black!55}, numbersep=7pt,
%  backgroundcolor=\color{ufvcodebg}, frame=single, rulecolor=\color{black!25},
%  breaklines=true, postbreak={\mbox{\textcolor{ufvcom}{$\hookrightarrow$}\space}},
%  extendedchars=true,
%  literate={á}{{\'a}}1 {é}{{\'e}}1 {í}{{\'i}}1 {ó}{{\'o}}1 {ú}{{\'u}}1
%           {Á}{{\'A}}1 {É}{{\'E}}1 {Í}{{\'I}}1 {Ó}{{\'O}}1 {Ú}{{\'U}}1
%           {â}{{\^a}}1 {ê}{{\^e}}1 {ô}{{\^o}}1 {ã}{{\~a}}1 {õ}{{\~o}}1
%           {Ã}{{\~A}}1 {Õ}{{\~O}}1 {ç}{{\c c}}1 {Ç}{{\c C}}1 {à}{{\`a}}1
%           {º}{{\textordmasculine}}1 {ª}{{\textordfeminine}}1,
%}
% versão em preto e branco do mesmo estilo (para impressão sem cores)
%\lstdefinestyle{ufvcodigopb}{style=ufvcodigo, keywordstyle=\bfseries,
%  commentstyle=\itshape, stringstyle=\relax, backgroundcolor=\color{white},
%  rulecolor=\color{black}, numberstyle=\tiny}
%\lstset{style=ufvcodigo}

% ---
% REFERÊNCIAS: biblatex com o estilo abnt (ajustes UFV no abntex2-UFV.sty)
% ---
\usepackage[backend=biber,
            style=abnt,
            uniquename=init,
            giveninits,
            repeatfields,
            justify,
            noslsn]{biblatex}

\usepackage{xurl}               % quebra de linha em URLs longas
\urlstyle{same}

\addbibresource{referencias.bib}
\setlength\bibparsep{10pt}
\setlength{\emergencystretch}{3em}

% ---
% SIGLAS (pacote acro, carregado pelo abntex2-UFV)
% Declare cada sigla aqui e use \ac{chave} na primeira ocorrência no texto (a sigla sai por
% extenso, seguida da sigla entre parênteses) e nas seguintes (só a sigla). Não use \ac em
% resumos, títulos e legendas. Nomes em outro idioma ficam em itálico (\textit{...}).
% ---
%\DeclareAcronym{abnt}{short = ABNT, long = Associação Brasileira de Normas Técnicas}
%\DeclareAcronym{api}{short = API, long = Interface de Programação de Aplicações}
% Abreviaturas (lista própria): use tag = abreviatura e \acs{chave}
%\DeclareAcronym{art}{short = art., long = artigo, tag = abreviatura}

% ---
% DADOS DO TRABALHO (capa, folha de rosto e folha de aprovação)
% ---
\titulo{Título do trabalho}
\autor{Nome do Autor}
\local{Rio Paranaíba}
\data{\the\year}
\orientador{Nome do Orientador}      % \orientadora{...} no feminino; instituição padrão: UFV-CRP
%\coorientador{Nome do Coorientador} % \coorientadora{...} no feminino
\instituicao{Universidade Federal de Viçosa}

\campus{\emph{Campus} Rio Paranaíba}   % aparece na capa só com a opção capacampus
\curso{Sistemas de Informação}
\membrobanca{Dr. Nome do Membro da Banca}       % instituição padrão: UFV-CRP
\membrobanca[UFV]{Dra. Nome do Membro da Banca} % ou informe a instituição entre colchetes
%\databanca{14 de janeiro de 2026}              % se omitida, a folha de aprovação traz espaço em branco

% Natureza do trabalho (nota descritiva da folha de rosto), conforme o manual da UFV:
% \preambuloUFV{tcc|bacharelado|licenciatura}{título}   ou   \preambulo{texto livre}
\preambuloUFV{bacharelado}{Sistemas de Informação}

% ---
% Configurações de aparência do PDF final (para impressão, troque as cores dos links por black)
\makeatletter
\hypersetup{
        pdftitle={\@title},
        pdfauthor={\@author},
        pdfsubject={\imprimirpreambulo},
        pdfcreator={LaTeX with abnTeX2},
        colorlinks=true,
        linkcolor=black,
        citecolor=blue,
        filecolor=magenta,
        urlcolor=blue,
        bookmarksdepth=4
}
\makeatother
% ---

\begin{document}
\frenchspacing

% ----------------------------------------------------------
% ELEMENTOS PRÉ-TEXTUAIS
% ----------------------------------------------------------
\pretextual

\imprimircapa

% Folha de rosto (a ficha catalográfica só é obrigatória em teses e dissertações)
\imprimirfolhaderosto

\imprimirfolhadeaprovacao

% Dedicatória (opcional; sem título, texto na parte inferior à direita)
%\begin{dedicatoria}
%   Texto da dedicatória
%\end{dedicatoria}

% Agradecimentos (opcional em TCC)
\begin{agradecimentos}
% Escreva os agradecimentos aqui.
\end{agradecimentos}

% Epígrafe (opcional; sem título, texto na parte inferior à direita)
%\begin{epigrafe}
%    ``Texto da epígrafe.''\\
%    (Autor)
%\end{epigrafe}

% RESUMO em português (obrigatório): parágrafo único, 150 a 500 palavras, verbo na voz
% ativa e na terceira pessoa, com objetivo, método, resultados e conclusões.
\begin{resumo}
 \noindent
 % Escreva o resumo aqui.

 \palavraschave{palavra um, palavra dois, palavra três}
\end{resumo}

% RESUMO em inglês (obrigatório)
\begin{resumo}[Abstract]
\begin{otherlanguage*}{english}
   \noindent
   % Write the abstract here.

   \keywords{keyword one, keyword two, keyword three}
 \end{otherlanguage*}
\end{resumo}

% Listas de ilustrações: uma lista própria para cada tipo. Descomente as dos tipos que
% existem no seu trabalho (e remova as demais).
%\pdfbookmark[0]{\listfigurename}{lof}
%\listoffigures*
%\cleardoublepage

%\pdfbookmark[0]{\listofquadrosname}{loq}
%\listofquadros*
%\cleardoublepage

%\pdfbookmark[0]{\listofgraficosname}{lgr}
%\listofgraficos*
%\cleardoublepage

%\pdfbookmark[0]{\listofcodigosname}{loc}
%\listofcodigos*
%\cleardoublepage

%\pdfbookmark[0]{\listtablename}{lot}
%\listoftables*
%\cleardoublepage

% Lista de siglas e lista de abreviaturas (opcionais): geradas pelo acro com as siglas e
% abreviaturas declaradas e usadas no texto.
%\listadesiglas
%\listadeabreviaturas

% Lista de símbolos (opcional)
%\begin{simbolos}
%\item[$\infty$ ] Infinito
%\end{simbolos}

% Sumário
\pdfbookmark[0]{\contentsname}{toc}
\tableofcontents*
\cleardoublepage

% ----------------------------------------------------------
% ELEMENTOS TEXTUAIS
% ----------------------------------------------------------
\textual

% Cada seção primária fica em um arquivo da pasta capitulos/.
% Adicione, remova ou renomeie conforme a estrutura do seu trabalho.
% Observação: no abntex2, \chapter corresponde à "seção primária" da NBR 6024. O manual da UFV
% pede que o trabalho seja organizado em seções: o título impresso não leva a palavra
% "Capítulo", apenas o número e o título em caixa alta.
\chapter{Introdução}\label{sec:introducao}

% Contexto, problema, objetivo, abordagem, contribuição e organização do texto.

\section{Objetivos}\label{sec:objetivos}

\subsection{Objetivo geral}

% Um único objetivo, com verbo no infinitivo.

\subsection{Objetivos específicos}

\begin{alineas}
  \item primeiro objetivo específico;
  \item segundo objetivo específico.
\end{alineas}

\section{Justificativa}\label{sec:justificativa}

\section{Organização do trabalho}\label{sec:organizacao}

\chapter{Referencial Teórico}\label{sec:RefTeorico}

% Conceitos necessários ao trabalho, organizados por temas, com síntese crítica.

\section{Fundamentação}\label{sec:fundamentacao}

\section{Trabalhos Relacionados}\label{sec:TrabRel}

\chapter{Metodologia}\label{sec:metodos}

% Classificação da pesquisa, etapas, instrumentos e materiais, critérios de avaliação,
% aspectos éticos, reprodutibilidade e limitações.

\chapter{Resultados e Discussão}\label{sec:resultados}

% Resultados na ordem dos objetivos específicos e discussão (interpretação, comparação com
% a literatura, limitações).

\chapter{Conclusão}\label{sec:conclusao}

% Conclusões correspondentes aos objetivos, contribuições, limitações e trabalhos futuros.

%% Use apenas no projeto (TCC I) e ative o \input correspondente no arquivo principal.
% Três formas de apresentar o cronograma (quadro com marcas, quadro com células
% sombreadas e diagrama de Gantt) estão no guia, em docs/capitulos/06-cronograma.tex.
\chapter{Cronograma}\label{sec:cronograma}
   % apenas no projeto (TCC I)

% ----------------------------------------------------------
% ELEMENTOS PÓS-TEXTUAIS
% ----------------------------------------------------------
\postextual

% Lista de referências (o aviso "Empty bibliography" é normal enquanto não houver citações)
\printbibliography

% Glossário (opcional): relação de termos técnicos com suas definições
%\begin{glossario}
%\termo{Termo}{definição do termo.}
%\end{glossario}

% Apêndices (texto elaborado pelo autor)
%\begin{apendicesenv}
%\partapendices
%\chapter{Primeiro apêndice}
%\end{apendicesenv}

% Anexos (documentos de terceiros)
%\begin{anexosenv}
%\partanexos
%\chapter{Primeiro anexo}
%\end{anexosenv}

\end{document}
